\section{Introduction}

Cyber-physical systems (CPS) tightly integrate computation, communication, and physical processes. Industrial PLCs continuously exchange actuator commands, sensor readings, and setpoints with connected machines over field protocols such as Siemens S7, Modbus, and OPC-UA. This control-plane communication stream is the operational nervous system of an industrial plant: every polling cycle carries the commands that move valves, start conveyors, and adjust setpoints.

\textbf{Problem.} An adversary with protocol-level network access to a PLC can read this stream, infer the current process state, and inject writes timed to cause physical damage — overflowing a tank, missort items on a conveyor, or triggering a safety alarm. A defender watching the same stream can detect writes that are semantically inconsistent with the observed controller behavior and block them before they execute. Yet most LLM-for-CPS security studies remain offline or simulation-only. For example, AttackLLM~\cite{ahmed2025attackllm} evaluates generated attacks on static traces, and autonomous cyber agents such as PentestGPT~\cite{deng2024pentestgpt} and Hackphyr~\cite{rigaki2024hackphyr} focus on IT infrastructure rather than physics-constrained communication streams. Adaptive frameworks such as Evo-Defender~\cite{hu2025evodefender} study attacker--defender co-evolution but do not incorporate LLM reasoning, while MALF~\cite{ning2025malf} operates at protocol level rather than process level. As a result, it remains unclear whether LLMs can perform effective security analysis of live control-plane communication --- where decisions must account for real controller timing, process phase, and the physical consequences of each write.

\textbf{Motivation and insight.} The live communication stream between PLC and plant carries exactly the information needed for both attack and defense reasoning: current sensor values, actuator states, setpoints, and control trends. Our key insight is that by placing a local LLM as a man-in-the-middle agent in this stream --- reading the same tags an attacker or defender would observe --- we can automate both roles systematically and symmetrically on real hardware. The same model, same observation pipeline, and same context structure serve as both red team and blue team, differing only in the prompt and the action they take.

\textbf{How we solve it.} We propose \textbf{CPSForge}, a framework that places a local LLM as a man-in-the-middle agent in real PLC communication, automating both red team (attack injection) and blue team (write interception) on the live tag stream. Every 500\,ms, CPSForge reads the PLC tag stream via S7 protocol (python-snap7). Every three polling cycles, it assembles a context-tier-specific prompt from the observed communication data and queries the LLM: as red team, decide to \emph{attack} (inject a write) or \emph{wait}; as blue team, decide to \emph{allow} or \emph{block} a proposed write. All writes are mediated by a mandatory safety shield before reaching the PLC. The framework is configuration-driven: the scene YAML provides the semantic interpretation of the tag stream (tag descriptions, valid ranges, attack surface, control objectives) for each plant. Connecting a different plant --- water treatment, power grid substation, manufacturing line --- requires only a new scene YAML; no code changes are required. We evaluate CPSForge on three Factory~I/O scenes connected to a real Siemens S7-1200 PLC.

Accordingly, the primary scientific contributions of this work are:

\begin{itemize}[leftmargin=1.5em]
    \item \textbf{C1.~Closed-loop LLM red team on live PLC communication.} We introduce an online LLM agent that operates as a man-in-the-middle in the PLC control loop: observe plant state $\to$ infer process phase $\to$ build context $\to$ decide (attack$|$wait) $\to$ act $\to$ observe result. Prior LLM-based CPS attack work either generates attacks offline without executing them (AttackLLM~\cite{ahmed2025attackllm}), operates at the network-packet layer rather than process-level tags (MALF~\cite{ning2025malf}), or targets IT infrastructure rather than physics-constrained control streams (PentestGPT~\cite{deng2024pentestgpt}). HARVEY~\cite{garcia2017harvey} demonstrated closed-loop PLC MitM but required a hand-crafted physics model per plant. CPSForge replaces plant-specific models with general-purpose LLM reasoning plus a configuration-driven scene abstraction, enabling the same agent to attack different plants without code changes.
    \item \textbf{C2.~Context-ablation study linking adversary knowledge to attack capability.} We define three strictly controlled context tiers---minimal (tag values only), partial (values + history + descriptions), and full (+ phase inference + prior action effects + attack guidance)---that map to realistic adversary reconnaissance levels: passive eavesdropping, field-device enumeration, and full plant documentation access. This enables the first controlled measurement of the \emph{minimum} process knowledge an LLM adversary needs to mount effective CPS attacks---a question that prior LLM-for-CPS work has not isolated because it evaluates only a single (typically maximal) context configuration.
    \item \textbf{C3.~Dual-objective fine-tuning study.} We design a QLoRA fine-tuning pipeline that trains Qwen2.5-3B-Instruct on 2{,}057~examples covering both verified successful attack traces and defense decisions (block/allow). Unlike Hackphyr~\cite{rigaki2024hackphyr}, which fine-tunes for IT penetration testing only, our dual-objective training enables a controlled before/after comparison of how task-specific CPS training changes \emph{both} red team and blue team capability within the same model---testing whether attack-aware fine-tuning transfers to defense.
    \item \textbf{C4.~Symmetric LLM blue team with defense comparison.} We propose an LLM-based blue team agent that uses the \emph{same model architecture, context pipeline, and inference infrastructure} as the red team, creating a symmetric attacker--defender pairing on the live tag stream. Unlike Evo-Defender~\cite{hu2025evodefender}, which co-evolves non-LLM strategies on simulators, and LLM-based IDS approaches~\cite{ann2024ids_llm, sharshar2025explainable} that classify historical traffic logs, our LLM defender evaluates each proposed PLC write in real time before execution. We compare seven defense configurations---no defense, rule-based (threshold, phase-aware, intent-consistency), LLM-only, and combined---enabling the first systematic comparison of rule-based vs.\ LLM-based vs.\ hybrid runtime defense against LLM-generated CPS attacks.
    \item \textbf{C5.~Cross-model robustness evaluation.} We evaluate three local models from two families---Qwen2.5-3B-Instruct, Qwen3-1.7B, and SmolLM3-3B---to test whether observed attack and defense trends are robust across architectures or strongly model-specific.
    \item \textbf{C6.~Frontier model scaling study.} We extend the evaluation with GPT-4o-mini (a frontier API model) across the same three scenes, demonstrating that self-deterrence vanishes at frontier scale, ASR reaches 100\% on Sorting by Weight, and the local LLM defender generalizes cross-model to block frontier-model attacks.
    \item \textbf{C7.~Comprehensive PLC-in-the-loop evaluation.} We evaluate C1--C6 through 333~live experiment runs across three Factory~I/O scenes on a real Siemens S7-1200 PLC, with $\geq$3~repeats per cell and unified per-step logging for reproducibility.
\end{itemize}

The remainder of the paper is organized as follows. Section~\ref{sec:background} provides background and motivation. Section~\ref{sec:threat} defines the threat model and assumptions. Section~\ref{sec:design} presents the CPSForge architecture. Section~\ref{sec:implementation} describes the implementation over a real PLC and Factory~I/O. Section~\ref{sec:methodology} describes the evaluation methodology. Section~\ref{sec:evaluation} presents experimental results. Section~\ref{sec:discussion} discusses findings and limitations. Section~\ref{sec:related} reviews related work, and Section~\ref{sec:conclusion} concludes.
